
\tphe{} is compared here against the literature it draws on and is answerable to, not against a
systematic review's full search strategy. Unless flagged otherwise, every citation below offers
adjacent, analogous, or precedent evidence for a construct \tphe{} uses --- never evidence that
\tphe{} itself, or any of its five rules, works. The full ten-topic narrative this synthesis
condenses, including the burden-of-disease case, the preconception-medicine analogue, the
peer-support and technology-facilitated-abuse literatures, the assessor-independence analogy,
decision-quality and digital self-assessment tool literature, language-access evidence, and the
tier-one Islamic bioethics companion literature, together with the evidence map for propositions
P1--P7 (Table~\ref{tab:evidencemap}), is reproduced in full, with every citation intact, in
Appendix~E (Extended evidence review).

\textbf{\tphe{} is proposed as an addition, not a replacement, to structures that already exist in
this region.} WHO's RESPECT framework already names relationship-skills strengthening as one of
seven population-level violence-prevention strategies, though a global review of reviews updating
it found only 6 of 178 eligible reviews targeted that strategy specifically \cite{who_respect,ref30}
--- a strategy named, not yet built out. Thailand's own Health Region~12 premarital course and
Malaysia's administratively mandatory Kursus Pra-Perkahwinan Islam, examined here only through a
moral-education content analysis at one Kuala Lumpur centre \cite{charoenwong2025}, show that a
premarital contact point can already be delivered at population scale in Muslim and Thai settings.
Neither is shown, in the sources available to this review, to carry a stop rule, a private channel
independent of the couple, or a pre-specified failure criterion. \tphe{} is offered as a governance
layer that could be added to structures such as these, not as a competing curriculum (full
comparison: Table~\ref{tab:novelty}, \S4.2).

\textbf{Knowledge gained at a premarital contact point does not reliably convert into the outcome
that contact point exists for.} Premarital genetic screening is the closest direct precedent:
premarital thalassaemia education raises knowledge ($p<0.001$) without reliably raising uptake
\cite{ref13}, and premarital sickle-cell-trait screening uptake across Africa pools under half of
eligible couples (47.82\%, 95\% CI 46.53--49.11), described as ``suboptimal'' \cite{ref12}. The
relationship-education literature shows the same dissociation at larger scale: across cumulative
meta-analyses, effects on relationship quality and communication skill are consistently small
($d\approx0.03$--$0.45$) \cite{ref6,ref7}; the most current synthesis of a government-funded
programme (Hawkins, 2022) found small effects on relationship quality ($d=0.114$) and skills
($d=0.132$) and non-significant effects on stability, parenting, and child well-being \cite{ref7};
and a meta-analysis of relationship aggression specifically (Karantzas, 2023) found $d\approx0.11$
overall, not statistically convincing when restricted to RCTs alone ($d\approx0.05$) \cite{ref8}.
None of this literature has measured informed choice, coercion recognition, or safe help-seeking.

\textbf{The same identification-without-response pattern motivates R4's referral-at-connection
rule.} A Cochrane review of 13 RCTs found IPV screening increased identification (OR 2.95) but no
evidence of an effect on referral (OR 2.24, 2 studies only) or on re-exposure at 3--18 months
\cite{ref38}. The 2025 USPSTF evidence report reviewed 35 studies ($n=18{,}358$): 3 RCTs
($n=3{,}759$) comparing screening with no screening found no significant reduction in IPV or other
outcomes over 3--18 months \cite{ref39}, and the companion 2025 recommendation statement retained
only a moderate-net-benefit judgement for reproductive-age women \cite{ref40}. Screening finds; it
does not, by itself, connect.

\textbf{Recognition and behaviour dissociate in the literature closest to Layer~0's premise.} The
strongest cautionary evidence for a proposed extension built around raising recognition is a 1-year
cluster RCT of Coaching Boys Into Men (16 high schools, $n=1{,}513$ male athletes): perpetration
fell, but recognition of abusive behaviours did not change (95\% CI $-0.15$ to $0.09$), and
bystander intentions were unchanged \cite{ref190} --- recognition and behaviour separating exactly
as \tphe{}'s failure criterion is built to catch. Layer~0's proposed anonymous self-check has a
direct, cautionary precedent of its own: I-DECIDE, a masked two-arm RCT ($n=422$, 79--80\% 12-month
retention), found no improvement in self-efficacy at 6 or 12 months, with scores slightly lower in
the intervention arm than control (mean difference 1.3, 95\% CI 0.3--2.3 at 6 months) and no
between-group difference in depressive symptoms \cite{ref119,ref120}. That null result
is reported here, not set aside, because it is the closest outcome-domain precedent this proposed
extension has.

\textbf{The stop rule's own evidence base is small and explicitly cautious.} The only systematic
review of couples therapy as an IPV treatment, after screening 1{,}733 citations, found only 6
studies met inclusion; pooled data suggested conjoint treatment ``viable in select situations,''
with professionals cautious about violent-retaliation risk \cite{ref45}. Survivor-voiced evidence
complicates any default toward joint response: only 9 of 17 women in one qualitative sample rated
couples counselling a good idea \cite{ref48}. This literature supports withholding routine joint
work where coercive control is confirmed; it does not demonstrate that any stop rule, including
\tphe{}'s, improves outcomes.

\textbf{Cultural-competence training repeatedly raises provider-level knowledge or attitudes without
a corresponding patient-outcome gain} --- the pattern R1's cultural-intelligibility safeguard is
built to answer, not to assume already solved. A Cochrane review of 5 RCTs (337 professionals,
8{,}400 patients) rated the evidence low-quality, with some improvement in involvement or
understanding but no reliable clinical-outcome effect \cite{ref55}; a narrower review found no
significant patient-outcome improvement in any of the studies it identified \cite{ref56}; and
Indigenous cultural-safety training across four high-income countries found outcomes overwhelmingly
provider-level, with weak clinical-outcome evidence \cite{ref57}.

\textbf{Religiosity is not a uniform protective factor, which is why the ethical frame (\S3.3) is
presented as a declared duty rather than as evidence.} A 67-study systematic review and
meta-analysis found religiosity and spirituality protective for general interpersonal aggression but
\emph{not} specifically for domestic violence \cite{ref220}, complicating any simple ``faith
protects'' reading and reinforcing why Table~\ref{tab:ethic} pairs each ethical term with a
separately evidenced construct rather than treating the term itself as a finding.

\textbf{Direct Thai evidence remains sparse and geographically uneven, and none of it evaluates a
premarital programme.} A four-region national survey ($n=2{,}462$ married or cohabiting women) found
roughly 15\% lifetime psychological, physical, or sexual IPV \cite{ref72}; among 700 Muslim
adolescents on the southern border, 9\% were sexually experienced, with condom non-use, STI, and
repeat pregnancy concentrated in that subgroup \cite{ref76}; and among 360 Muslim army conscripts in
the three deep southern provinces, 78.7\% showed poor HIV/STI-transmission knowledge \cite{ref77}.
No study identified in this review evaluates a premarital-education programme, of any design, with a
Thai Muslim or Thai intercultural population.

\textbf{Read together, this literature does two things and stops short of a third.} It makes a
premarital contact point plausible to build and deliver at population scale (preconception medicine,
premarital genetic screening, and Thailand's and Malaysia's own existing courses); it repeatedly
shows that knowledge, attendance, satisfaction, and provider-level attitude gains do not reliably
convert into the behavioural, referral, or patient-outcome change a programme is built to produce
(screening, relationship education, and cultural-competence training alike). It does not show that
any premarital governance layer, of this or any design, changes informed choice, coercion
recognition, or safe help-seeking --- because no such test has yet been run. \tphe{} is offered as a
falsifiable synthesis awaiting that test, never as evidence that premarital education of this design
improves, prevents, or reduces any outcome. The full evidence base this synthesis draws from,
organised by proposition (P1--P7) and by topic, is in Appendix~E.
